\section{Application II: Polylog–Zeta Factorization and Variants}\label{sec:application-key-results}

This section presents further results that demonstrate the flexibility and depth of the FD-Lift framework. These include an exact factorization of the spectrum of the prime indicator function $\mathfrak{F}(z,q)$ and the construction of analytic $q$-analogs for other classical arithmetic functions.

\subsection{The Polylog-Zeta Factorization}

A direct application of the spectral property of the FD-Lift (Theorem~\ref{thm:fd-lift-core}) to the sum part $\mathfrak{S}(n,q)$ of the prime indicator function yields a factorization involving the Riemann Zeta function and the polylogarithm.

\begin{theorem}[Polylog-Zeta Factorization]\label{thm:polylog-zeta}
For $|q|>1$ and $\Re s>1$, the Dirichlet series for $\mathfrak{S}(n,q)$ and $\mathfrak{F}(n,q)$ are given by
\begin{align*}
\sum_{n\ge2}\frac{\mathfrak{S}(n,q)}{n^{s}}
&=(q-1)q\,\zeta(s)\,\bigl(\operatorname{Li}_{s}(1/q)-q^{-1}\bigr), \\
\sum_{n\ge2}\frac{\mathfrak{F}(n,q)}{n^{s}}
&=(q-1)q\,\bigl(\zeta(s)-1\bigr)\,\bigl(\operatorname{Li}_{s}(1/q)-q^{-1}\bigr).
\end{align*}
\end{theorem}
\begin{proof}[Sketch]
Use $\mathfrak S(n,q)=(q\!-\!1)q\sum_{d\mid n,\, d\ge2}q^{-d}$, swap sums to get $\zeta(s)\sum_{d\ge2}q^{-d}d^{-s}=\zeta(s)\big(\mathrm{Li}_s(1/q)-q^{-1}\big)$.
For $\mathfrak F$ subtract in addition the Dirichlet series of the corrector, $\sum_{n\ge2}(q-1)q\,q^{-n}n^{-s}=(q-1)q\big(\mathrm{Li}_s(1/q)-q^{-1}\big)$, which replaces $\zeta(s)$ by $\zeta(s)-1$.
Full derivation: Appendix~\ref{app:fdlift-proofs}.
\end{proof}

\begin{corollary}[Reconstruction of $\zeta(s)$]\label{cor:reconstruct-zeta}
For $|q|>1$ and $\Re s>1$ with $\operatorname{Li}_s(1/q)\neq q^{-1}$, the quotient
\[
\zeta(s)=\frac{\displaystyle \sum_{n\ge2}\mathfrak{S}(n,q)\,n^{-s}}{(q-1)q\big(\operatorname{Li}_s(1/q)-q^{-1}\big)}
\]
is well-defined. The condition $\operatorname{Li}_s(1/q)\neq q^{-1}$ holds in particular for real $q>1$ and real $s>1$, and for all $|q|\ge2$ and $\Re s>1$.
\end{corollary}
\begin{remark}
For real $q>1$ and real $s>1$ one has $\operatorname{Li}_s(1/q)-q^{-1}=\sum_{n\ge2}q^{-n}n^{-s}>0$. For $|q|\ge2$ and $\sigma=\Re s>1$, with $w=1/q$,
\[
\Bigl|\sum_{n\ge2}w^{n}n^{-s}\Bigr|\ \ge\ |w|^2\,2^{-\sigma}-\frac{|w|^3\,3^{-\sigma}}{1-|w|}\ \ge\ |w|^2\bigl(2^{-\sigma}-3^{-\sigma}\bigr)\ >\ 0 .
\]
For $1<|q|<2$ the condition is not automatic, even for real $q$: for instance $\operatorname{Li}_s(1/q)=q^{-1}$ at $q=1.01$ and $s\approx1.380809+23.326651\,\mathrm{i}$, where the quotient is not defined.
\end{remark}

\subsection{Analytic q-Analogs of Arithmetic Functions}
The presented approach can be generalized to construct analytic q-analogs of other classical arithmetic functions~\cite{gasper2004}.

\subsubsection{Definition and analyticity of $\mathfrak{F}_\tau(z,q)$ and $\mathfrak{F}_\sigma(z,q)$}
By choosing different weighting functions and removing the normalization factor, entire functions can be defined.
\begin{definition}[q-Analog of the Divisor-Counting Function]
\[\mathfrak{F}_\tau(z,q) = \left( \sum_{i=2}^{\infty} \frac{F(z,i)}{i^2} \left(\frac{1}{q}\right)^i \right) - \left(\frac{1}{q}\right)^z\]
\end{definition}
\begin{definition}[q-Analog of the Sum-of-Divisors Function]
\[\mathfrak{F}_\sigma(z,q) = \left( \sum_{i=2}^{\infty} \frac{F(z,i)}{i} \left(\frac{1}{q}\right)^i \right) - z \left(\frac{1}{q}\right)^z\]
\end{definition}
The analyticity of $\mathfrak{F}_\tau(z,q)$ follows directly from the proof for $\mathfrak{F}(z,q)$, as the summation terms are identical. Replacing the corrector $q^{-z}$ by the tangent-matched corrector of Section~\ref{sec:periodic-normalizer} gives the improved q-analog
\[
\mathfrak{F}_\tau^\sharp(z,q):=\sum_{i=2}^{\infty}\frac{F(z,i)}{i^2}\,q^{-i}-q^{-z}\bigl(1+(\log q)\,S_1(z)\bigr)=\frac{\Fsharp(z,q)}{(q-1)q},
\]
which has the same integer values as $\mathfrak F_\tau$ and the same real zeros as $\Fsharp$.

\subsubsection{Theorem (Connection to classical functions)}
\begin{theorem}\label{thm:q-analog-limit}
For any integer $n \ge 2$, the q-analogs generate divisor polynomials.
In the limit $q \to 1^+$, they recover the classical arithmetic functions:
\begin{align*}
    \mathfrak{F}_\tau(n,q) &= \sum_{i|n, 2 \le i < n} \left(\frac{1}{q}\right)^i & \xrightarrow{q\to1^+} \quad \tau(n)-2 \\
    \mathfrak{F}_\sigma(n,q) &= \sum_{i|n, 2 \le i < n} i \left(\frac{1}{q}\right)^i & \xrightarrow{q\to1^+} \quad \sigma(n)-n-1
\end{align*}
\end{theorem}

\noindent
For $z\notin\N$ the limit $q\to1^+$ does not define the classical functions, since the infinite series do not converge in that limit outside integer arguments.

\subsection{Interpretation as divisor polynomials}
Instead of a single numerical value, these functions yield a polynomial in $1/q$ whose structure encodes arithmetic information about the argument. For an integer $n$, they generate a polynomial in $1/q$ whose exponents encode the divisors of $n$ and whose coefficients encode arithmetic weights. This yields a representation where the divisors are explicitly encoded as exponents.

The role of the parameter $q$ is visualized in Figure~\ref{fig:q_analog_interpolation}, which displays $\mathfrak{F}_\tau(z,q)$ for several values of $q$. In the limiting case $q=1$, the function exactly reproduces the values of the classical divisor-counting function $\tau(n)-2$ at integer arguments. As $q$ increases, the curves deviate more significantly from these classical values and become flatter, since the weighting terms $q^{-d}$ in the divisor polynomial decrease. The parameter $q$ thus acts as a \emph{deformation parameter}, bridging classical arithmetic and a family of analytic functions.

\begin{figure}[!htbp]
\centering
\includegraphics[width=0.95\textwidth]{plot_q_analog_interpolation.pdf}
\caption{Comparison of the analytic q-analog $\mathfrak{F}_\tau(z,q)$ for several values of $q$. The orange diamonds represent the classical divisor-counting function $\tau(n)-2$, which corresponds to the exact case $q=1.0$. The solid curves show the function for $q = 1.01, 1.1, 1.2, 1.5,$ and $2.0$. The plot clearly illustrates that as $q$ increases, the curves become flatter and deviate more from the discrete classical values.}
\label{fig:q_analog_interpolation}
\end{figure}

\FloatBarrier

% ---------------------------------------------------------------------------
% ---------------------------------------------------------------------------
\subsection{Euler-factor design and sample numerics}
The FD-Lift separates local and spectral data. Two compact examples illustrate the Euler-factor viewpoint and supply small-scale numerics.

\subsubsection*{Example 1: $a(n)=n^{-2}$}
Here $A(s)=\sum_{n\ge1}\frac{a(n)}{n^{s}}=\sum_{n\ge1}n^{-(s+2)}=\zeta(s+2)$, hence
\[
\sum_{n\ge1}\frac{\mathcal{T}_a(n)}{n^{s}}=\zeta(s)\zeta(s+2),\qquad
\mathcal{T}_a(n)=(a*1)(n)=\sum_{d\mid n} d^{-2}=\sigma_{-2}(n).
\]
By multiplicativity this yields, for $n=\prod p^{k}$,
\[
\mathcal{T}_a(n)
=\prod_{p^k\parallel n}\frac{1-p^{-2(k+1)}}{1-p^{-2}}.
\]
A compact prime-power table (sufficient to reconstruct all values by multiplicativity) is given by:
\begin{table}[h]
\centering
\caption{Local factors $\mathcal{T}_a(p^k)$ for $a(n)=n^{-2}$ at small prime powers.}
\begin{tabular}{r c}
\toprule
$p^k$ & $\mathcal{T}_a(p^k)$\\
\midrule
$2$   & \(\frac{5}{4}\) \\
$2^2$ & \(\frac{21}{16}\) \\
$3$   & \(\frac{10}{9}\) \\
$3^2$ & \(\frac{91}{81}\) \\
$5$   & \(\frac{26}{25}\) \\
$5^2$ & \(\frac{651}{625}\) \\
\bottomrule
\end{tabular}
\end{table}


\subsubsection*{Example 2: $a(n)=\chi_4(n)\,n^{-2}$}
Let $\chi_4$ be the primitive character modulo $4$. Then $A(s)=\sum_{n\ge1}\chi_4(n)n^{-(s+2)}=L(s+2,\chi_4)$ and
\[
\sum_{n\ge1}\frac{\mathcal{T}_a(n)}{n^{s}}=\zeta(s)\,L(s+2,\chi_4),\qquad
\mathcal{T}_a(n)=\sum_{d\mid n}\frac{\chi_4(d)}{d^{2}}.
\]
By multiplicativity, and since $\chi_4(2)=0$, one has the local factorization
\[
\mathcal{T}_a(n)=\prod_{\substack{p^k\parallel n\\ p\ \mathrm{odd}}}\frac{1-(\chi_4(p)/p^{2})^{k+1}}{1-\chi_4(p)/p^{2}},
\qquad\text{and the factor at }p=2\text{ equals }1.
\]
\emph{In particular, the contribution of the $2$-adic part vanishes: $\mathcal{T}_a(2^k m)=\mathcal{T}_a(m)$ for odd $m$.}
A compact comparative table for odd $n\le 11$ is:
\begin{table}[h]
\centering
\caption{Comparison of $\mathcal{T}_a(n)$ for $a(n)=n^{-2}$ and $a(n)=\chi_4(n)n^{-2}$ (odd $n\le 11$).}
\begin{tabular}{r c c}
\toprule
$n$ & $\mathcal{T}_{a=n^{-2}}(n)$ & $\mathcal{T}_{a=\chi_4(n)n^{-2}}(n)$\\
\midrule
1 & 1 & 1 \\
3 & \(\frac{10}{9}\) & \(\frac{8}{9}\) \\
5 & \(\frac{26}{25}\) & \(\frac{26}{25}\) \\
7 & \(\frac{50}{49}\) & \(\frac{48}{49}\) \\
9 & \(\frac{91}{81}\) & \(\frac{73}{81}\) \\
11 & \(\frac{122}{121}\) & \(\frac{120}{121}\) \\
\bottomrule
\end{tabular}
\end{table}


\noindent In particular,
\[
A(1)=\sum_{n\ge1}\frac{a(n)}{n}=
\begin{cases}
\zeta(3), & a(n)=n^{-2},\\[3pt]
L(3,\chi_4)=\beta(3)=\dfrac{\pi^{3}}{32}, & a(n)=\chi_4(n)n^{-2}.
\end{cases}
\]

% ---------------------------------------------------------------------------

